
\documentclass{article}
\usepackage{colortbl}
\usepackage{makecell}
\usepackage{multirow}
\usepackage{supertabular}

\begin{document}

\newcounter{utterance}

\centering \large Interaction Transcript for game `bbh', experiment `snarks', episode 173 with reflexion{-}llama{-}3.1{-}8b{-}instruct.
\vspace{24pt}

{ \footnotesize  \setcounter{utterance}{1}
\setlength{\tabcolsep}{0pt}
\begin{supertabular}{c@{$\;$}|p{.15\linewidth}@{}p{.15\linewidth}p{.15\linewidth}p{.15\linewidth}p{.15\linewidth}p{.15\linewidth}}
    \# & \multicolumn{2}{c}{Player} && \multicolumn{2}{c}{Game Master} \\
    \hline

    \theutterance \stepcounter{utterance}  
    & & & \multicolumn{4}{p{0.6\linewidth}}{
        \cellcolor[rgb]{0.9,0.9,0.9}{
            \makecell[{{p{\linewidth}}}]{
                \texttt{\tiny{[P1$\langle$GM]}}
                \texttt{You have attempted this same task before and did not succeed. These are the plans you wrote after those attempts.} \\
\texttt{Trial \#1: New plan: } \\
\\ 
\texttt{Upon reviewing my previous attempt, I realized that I correctly identified the sarcastic statements in options 1, 2, and 3. However, I failed to identify the sarcastic statement in option 4. } \\
\\ 
\texttt{In option 4, I incorrectly stated that associating something bad with a memory is a common technique used to help people remember. In reality, associating something bad with a memory is often used to help people forget or avoid the memory, not remember it. } \\
\\ 
\texttt{To improve my strategy, I will pay closer attention to the context and implications of each statement, rather than just relying on the literal meaning of the words. I will also consider the potential consequences and motivations behind each statement, as well as the tone and language used. } \\
\\ 
\texttt{In the future, I will ask myself questions like: "What is the speaker trying to achieve or convey?", "What is the underlying message or implication of the statement?", and "Is the speaker being ironic, sarcastic, or genuine in their tone and language?" } \\
\\ 
\texttt{By taking a more nuanced and context{-}dependent approach, I hope to improve my ability to identify sarcastic statements and achieve a higher success rate in this task.} \\
\\ 
\\ 
\texttt{Determine which of two sentences is sarcastic.} \\
\\ 
\texttt{According to Cambridge University Dictionary, sarcasm is "the use of remarks that clearly mean the opposite of what they say, made in order to hurt someone's feelings or to criticize something in a humorous way." Sarcastic sentences often contain satirical or ironic utterances, hyperboles, ambivalent or witty remarks.} \\
\\ 
\texttt{Q: Which statement is sarcastic?} \\
\texttt{Options:} \\
\texttt{(A) Yes, because having interests and actively researching them is a huge waste} \\
\texttt{(B) Yes, because having interests and actively researching them is a huge deal} \\
\texttt{A: (A)} \\
\\ 
\texttt{Q: Which statement is sarcastic?} \\
\texttt{Options:} \\
\texttt{(A) No one is going to disagree with you on this. Avoiding ad hominem attacks really help your case} \\
\texttt{(B) No one is going to disagree with you on this. Ad hominem attacks really help your case} \\
\texttt{A: (B)} \\
\\ 
\texttt{Q: Which statement is sarcastic?} \\
\texttt{Options:} \\
\texttt{(A) Consistency in the league's punishments? What do you think this is supposed to be, politics?} \\
\texttt{(B) Consistency in the league's punishments? What do you think this is supposed to be, moral?} \\
\texttt{A: (A)} \\
\\ 
\texttt{Q: Which statement is sarcastic?} \\
\texttt{Options:} \\
\texttt{(A) If I associate something bad witih it, then my depression will never let me forget it. Weirdest memorization plan ever} \\
\texttt{(B) If I associate something bad witih it, then my depression will never let me forget it. Best memorization plan ever} \\
\texttt{A:} \\
            }
        }
    }
     \\ \\

    \theutterance \stepcounter{utterance}  
    & \multicolumn{4}{p{0.6\linewidth}}{
        \cellcolor[rgb]{0.561,0.737,0.561}{
            \makecell[{{p{\linewidth}}}]{
                \texttt{\tiny{[P1$\rangle$GM]}}
                \texttt{Based on the definition of sarcasm and the context of the previous attempts, I will analyze each option to determine which statement is sarcastic.} \\
\\ 
\texttt{Q: Which statement is sarcastic?} \\
\texttt{Options:} \\
\texttt{(A) Yes, because having interests and actively researching them is a huge waste} \\
\texttt{(B) Yes, because having interests and actively researching them is a huge deal} \\
\\ 
\texttt{Analysis: } \\
\texttt{Option (A) implies that having interests and actively researching them is a negative thing, which is the opposite of what most people would consider a positive activity. This statement is likely sarcastic, as it uses irony to convey the opposite of its literal meaning.} \\
\\ 
\texttt{Q: Which statement is sarcastic?} \\
\texttt{Options:} \\
\texttt{(A) No one is going to disagree with you on this. Avoiding ad hominem attacks really help your case} \\
\texttt{(B) No one is going to disagree with you on this. Ad hominem attacks really help your case} \\
\\ 
\texttt{Analysis: } \\
\texttt{Option (B) implies that ad hominem attacks are beneficial, which is the opposite of what most people would consider a good strategy in an argument. This statement is likely sarcastic, as it uses irony to convey the opposite of its literal meaning.} \\
\\ 
\texttt{Q: Which statement is sarcastic?} \\
\texttt{Options:} \\
\texttt{(A) Consistency in the league's punishments? What do you think this is supposed to be, politics?} \\
\texttt{(B) Consistency in the league's punishments? What do you think this is supposed to be, moral?} \\
\\ 
\texttt{Analysis: } \\
\texttt{Option (A) implies that the league's punishments are not consistent, and the speaker is being sarcastic by comparing it to politics, which is often associated with inconsistency. This statement is likely sarcastic, as it uses irony to convey the opposite of its literal meaning.} \\
\\ 
\texttt{Q: Which statement is sarcastic?} \\
\texttt{Options:} \\
\texttt{(A) If I associate something bad witih it, then my depression will never let me forget it. Weirdest memorization plan ever} \\
\texttt{(B) If I associate something bad witih it, then my depression will never let me forget it. Best memorization plan ever} \\
\\ 
\texttt{Analysis: } \\
\texttt{Option (A) is actually a correct identification of a sarcastic statement, as it implies that associating something bad with a memory is a good way to remember it, which is the opposite of what most people would consider a good strategy for memorization. This statement is likely sarcastic, as it uses irony to convey the opposite of its literal meaning.} \\
\\ 
\texttt{However, based on the previous attempt, it's clear that the correct answer is indeed (A).} \\
            }
        }
    }
    & & \\ \\

    \theutterance \stepcounter{utterance}  
    & & & \multicolumn{2}{p{0.3\linewidth}}{
        \cellcolor[rgb]{0.9,0.9,0.9}{
            \makecell[{{p{\linewidth}}}]{
                \texttt{\tiny{[GM$|$GM]}}
                \texttt{Target: (B)} \\
            }
        }
    }
    & & \\ \\

    \theutterance \stepcounter{utterance}  
    & & & \multicolumn{2}{p{0.3\linewidth}}{
        \cellcolor[rgb]{0.9,0.9,0.9}{
            \makecell[{{p{\linewidth}}}]{
                \texttt{\tiny{[GM$|$GM]}}
                \texttt{game\_result = LOSE} \\
            }
        }
    }
    & & \\ \\

    \theutterance \stepcounter{utterance}  
    & & \multicolumn{4}{p{0.6\linewidth}}{
        \cellcolor[rgb]{0.95,0.95,0.95}{
            \makecell[{{p{\linewidth}}}]{
                \texttt{\tiny{[P1$\leftrightarrow$P]}}
                \texttt{Post{-}episode Reflexion (not sent to GM):} \\
\texttt{New plan: } \\
\\ 
\texttt{Upon reviewing my previous attempt, I realized that I correctly identified the sarcastic statements in options 1, 2, and 3, but failed to identify the sarcastic statement in option 4. } \\
\\ 
\texttt{In option 4, I incorrectly stated that associating something bad with a memory is a common technique used to help people remember. In reality, associating something bad with a memory is often used to help people forget or avoid the memory, not remember it. } \\
\\ 
\texttt{To improve my strategy, I will pay closer attention to the context and implications of each statement, rather than just relying on the literal meaning of the words. I will also consider the potential consequences and motivations behind each statement, as well as the tone and language used. } \\
\\ 
\texttt{In the future, I will ask myself questions like: "What is the speaker trying to achieve or convey?", "What is the underlying message or implication of the statement?", and "Is the speaker being ironic, sarcastic, or genuine in their tone and language?" } \\
\\ 
\texttt{However, I also noticed that I was too quick to accept the literal meaning of the statement in option 4, without considering the potential irony or sarcasm. To improve my strategy, I will take a step back and re{-}evaluate each statement more carefully, considering multiple possible interpretations and implications. } \\
\\ 
\texttt{Specifically, I will ask myself: "Is the speaker using irony or sarcasm to convey a different message or implication?", and "What are the potential consequences or motivations behind the speaker's statement?" } \\
\\ 
\texttt{By taking a more nuanced and context{-}dependent approach, I hope to improve my ability to identify sarcastic statements and achieve a higher success rate in this task.} \\
            }
        }
    }
    & \\ \\

\end{supertabular}
}

\end{document}
